\chapter{Measure Theory}
\section{Measure and measurable sets}

\subsection{Set theory basics}
For the development of measure theory, Borel has introduced the following defintion:
\begin{definition}[Limes superior and limes inferior]
	Given $(A_n)_{n \geq 1}$ a sequence of subsets in $X$. We define:
	\begin{align*}
		\limsup_n A_n &:= \{ x \in X: x \in A_n \textbf{ for infinitely many } A_n \}\\
		\liminf_n A_n &:= \{ x \in X: \exists n_0(x) \in \mathbb{N}\forall n \geq n_0(x): x \in A_n \} \textbf{ x is in all but finitely many $A_n$}
	\end{align*}
	equivalently we see:
	\begin{align*}
		\limsup_n A_n &= \cap_{n = 1}^\infty \cup_{k = n}^\infty A_n\\ 
		\liminf_n A_n &= \cup_{n = 1}^\infty \cap_{k = n}^\infty A_n
	\end{align*}
	We define the convergence of sets as :
	\begin{align*}
		\limsup_n A_n = \liminf_n A_n := \lim_n A_n
	\end{align*}
\end{definition}
Note that the definition of $\liminf$ coincides with the definition of limes we are familiar with. However for a sequence of sets it always exists. The limes of a sequence of sets could be interpreted as: all elements that appear infinitely many times is in $A_n$ for $n$ large enough.


\subsection{($\sigma)$- Ring and $\sigma$-Algebra}
\begin{theorem}
	Given a power set of $X$ $\mathfrak{B}(X)$, $(\mathfrak{B}(X), \Delta, \cap)$ is a commutative ring with the Null $\emptyset$ and the identity $X$.
\end{theorem}
\begin{proof}
	There is an isomorphism between the power set $\mathfrak{B}(X)$ and the ring of characteristic function(for each $A \subset X$ we can define a homomorphism $\chi_A: X \rightarrow \{\bar{1}, \bar{0} \}$ to the ring $\{\bar{0}, \bar{1}\}$). Obviously it's a isomorphism.
\end{proof}

The algebraic structure works for the finite case. For developing the integration theory, finite combination is not enough, therefore we introduce the concept of $\sigma$- ring and $\sigma$- field.

\begin{definition}[$\sigma$-Ring and Field]
Given a set $X$
	\begin{itemize}
		\item Ring and field: A subring of the ring defined above is a \textbf{ring }over $X$. Such ring is an \textbf{Field} if it contains the $X$.  
		\item $\sigma$-Ring  and $\sigma$-field: A ring is a \textbf{$\sigma$-ring} if it's closed under \textbf{countable union}
	\end{itemize}
\end{definition}

\begin{theorem}
	$\mathfrak{R} \subset \mathfrak B(X)$ is a ring iff one of the following is true:
	\begin{enumerate}
		\item $\emptyset \in \mathfrak{R}, \forall A, B \in \mathfrak{R}: A \Delta B \in \mathfrak{R}, A\bigcap B \in \mathfrak{R}$
		\item $\dots\Delta\dots \bigcup \dots$
		\item $\dots\backslash\dots \bigcup \dots$
	\end{enumerate}
\end{theorem}


\subsection{Generator of the Borel-sets}
In the following theorem, we introduced an useful technique called \textit{Principle of Good sets}. It's often applied when we want to prove that a generator produces the same set under 2 different sets of operations.
\begin{theorem}
	$f: X \rightarrow Y$ a measurable map, $\mathfrak{G} \subset \mathfrak{B}(X)$ a generator of $\mathfrak{B}$. we have:
	\begin{align*}
		\sigma(f^{-1}(\mathfrak{G}) ) = f^{-1}(\sigma (\mathfrak{G}))
	\end{align*}
\end{theorem} 
\begin{proof}
\end{proof}

\subsection{Semiring}
In practice it's often easy to generate the $\sigma$-ring or $\sigma$-field with generators of simple structure. A class of this type is what we'll introduce: the \textbf{Semi-ring}. This concept was first introduced by Johann von Neumann.
\begin{definition}[Semi-ring]
 A set of sets $\mathfrak{H}$ is called semi-ring if it satisfies
 \begin{enumerate}
 	\item $\emptyset \in \mathfrak{H}$
 	\item $\forall A, B \in \mathfrak{H}: A \cap B \in \mathfrak{H}$
 	\item $\forall A, B \in \mathfrak{H}$ there exists disjoint $C_1, C_2, \dots C_n: A\slash B = \bigcup_{k = 1}^n C_k$
 \end{enumerate}
\end{definition}
$\quad$\\
By \textbf{H.Hahn} we know semi-ring can generate a ring:
\begin{theorem}
	Is $\mathfrak{H} \subset \mathfrak{B}(X)$ a semiring, then the disjoint union of the element in $\mathfrak{H}$ generates a ring:
	\begin{align*}
		\mathfrak{R} := \{ \bigcup_{I < \infty} A_i,\quad \textbf{$A_i\in \mathfrak{H}$ disjoint } \} 
	\end{align*}
\end{theorem}
\begin{proof}
\end{proof}


\subsection{Monotone class and Dynkin System}
In the last 2 subsections, we have introduced the ring-structure of the set of sets. Also, given a semi-ring, we can generate a ring with the operation of disjoint union. Since the main object of study is the $\sigma$-field, we need to find the operation, under which a ring becomes a $\sigma$-field. In this spirit, we'll introduce the concept of \textbf{Dynkin system}, the main result is the following theorem:

	"Given an intersection-stable generator $\mathfrak{G} \subset \mathfrak{B}(X)$, the Dynkin-system generated by $\mathfrak{G}$ equals to the $\sigma$-field generated by $\mathfrak{G}$."
	 
\subsubsection*{The monotone class}
\begin{definition}[Monotone Class]
	A system is a monotone class if it is closed under countable union, i.e. $\mathfrak{M} \subset \mathfrak{B}(X)$
	\begin{align*}
		\forall (A_i)_{i \in I} \in \mathfrak{M} \textbf{ monotone}: \lim_{i\rightarrow \infty} A_i \in \mathfrak{M} \}
	\end{align*}
\end{definition}
The importance of the monotone class lies in the fact that it's the necessary and sufficient condition for the underlying ring to generate a $\sigma$-ring.

\begin{theorem}
	The $\sigma$-ring generated by the $\mathfrak{R} \subset \mathfrak{B}(X)$ equals to the monotone class generated by $\mathfrak{R}$.
\end{theorem}
\begin{proof}
	The claim is $\mathfrak{M}(\mathfrak{R}) = \sigma(\mathfrak{R})$\\
	"$\subset$": trivial \\
	"$\supset$": again we apply \textit{Principle of the good sets}. Given $\mathfrak{M}(\mathfrak{R}):= \mathfrak{M}$. We define the sets in $X$ that are "good" under the ring operation(as $\mathfrak{D}$), and show that $\mathfrak{M} \subset \mathfrak{D} $.\par 
	Given $A \in \mathfrak{M}$, we define:
	\begin{align*}
		\mathfrak{D}(A):= \{ B \subset X: A \backslash B \in \mathfrak{M}, B \backslash A \in \mathfrak{M}, B \bigcup A \in \mathfrak{M} \}
	\end{align*}
	
It's easy to verify that $\mathfrak{D}(A)$ is a monotone class, since the ring operation upon $A$ would enclose all the sets in $\mathfrak{R}$, and $\mathfrak{M}$ is the smallest monotone class generated by $\mathfrak{R}$, it follows $\mathfrak{R}  \subset \mathfrak{M} \subset \mathfrak{D}(A)$.\par
Which means for any $A \in \mathfrak{M}$, to get a set in $\mathfrak{M}$ after performing the ring operation with a set $B$ in $X$, we can select $B$ as any set in $\mathfrak{M}$. This means $\mathfrak{M}$ is closed under the ring operation. 
\end{proof}
\begin{remark}[principle of the good sets]
\end{remark}


\begin{definition}[Dynkin system($\lambda$ System)]
$\mathfrak{D}\subset \mathfrak{B}(X)$ is a Dynkin system if
\begin{enumerate}
	\item $X \in \mathfrak{D}$
	\item $\forall A \in \mathfrak{D}: A^c \in \mathfrak{D}$
	\item $\forall (A_i)_{i \geq 1} \in \mathfrak{D} \textbf{ disjoint }: \bigcup_{i = 1}^\infty A_i \in \mathfrak{D}$
\end{enumerate} 
An equivalent definition is given by replacing 2. with "$\backslash$" and 3. with \textbf{monotone}
\end{definition}
\begin{remark}
The first definition corresponds the definition of measure. The $\sigma$-additivity of the measure is perfectly mirrored by the properties of the Dynkin-system.\par 
The Dykin-system differs from a $\sigma$-field in that it's closed under countable union of \textbf{disjoint} or \textbf{monotone} sets, where the $\sigma$-field of arbitrary sets.\par 
As mentioned at the beginning, the following most important theorem of this subchapter provides a tool to prove the uniqueness of the measure:
\end{remark}

\begin{lemma}
	A Dykin-system is a $\sigma$-field if it's intersection-stable.
\end{lemma}
\begin{proof}
	From sequence of arbitrary sets $A_n$ we may construct a sequence of monotone sets by:
	\begin{align*}
		B_1 &= A_1\\
		B_n: &= A_n \backslash B_{n-1}
	\end{align*}
	Then we have
	\begin{align*}
		\bigcup A_n = \bigcup B_n = \bigcup(A_n \cap B_{n-1}^C)
	\end{align*}
	If the Dykin-system is stable under intersection then we know it's stable under arbitrary union
\end{proof}

\begin{theorem}
	Given a intersection-stable generator $\mathfrak{G}$, the Dykin-system and $\sigma$-field generated by the generator is the same.
\end{theorem}







\subsection{Content, premeasure, and measure}
In the last section we dealt with the problem of the domain of a measure. From now on we shift our topic to the measure function. We follow the induction process of semi-ring, $\sigma$-ring and $\sigma$-algebra to introduce the measure. Axiomatically, a measure would satisfies positiveness and additivity.
\begin{definition}[Content, premeasure and measure]
Given $X$ a set, we define:
	\begin{itemize}
		\item \textbf{Content}: defined on a semiring: $\mathcal{Z} \rightarrow \mathbb{R}^+$ satisfies:
		\begin{enumerate}
			\item $\mu(\emptyset) = 0$. 
			\item $\mu > 0$
			\item \textbf{Finite additivity:} $A_i$ disjoint: $\mu(\bigcup_{i = 1}^n A_i) = \sum_{i = 1}^n \mu(A_i)$
		\end{enumerate}
		\item \textbf{premeasure}: function with properties 1), 2) whereas in 3) we would have $\sigma$-additivity.
		\item \textbf{measure}: defined on a $\sigma$-field.
	\end{itemize}
\end{definition}


\subsection{Extension from premeasure to measure}
Now we have a premeasure defined on a semi-ring. The tools at hand are 1) every set in the $\sigma$-field generated by $\mathfrak{H}$ could be approximated by the union of sets in $\mathfrak{H}$, 2) The $\sigma$-additivity, i.e.:
\begin{align*}
	\forall A \in \sigma(\mathfrak{R}): A = \bigcup_{n =1}^\infty A_n, \quad  A_n \in \mathfrak{R} \textbf{ disjoint }
\end{align*}
therefore
\begin{align*}
	\mu(A) = \sum_{n = 1}^\infty \mu(A_n)
\end{align*}
The problem is, the choice of $A_n$ is not unique. For an arbitrary $A \in \sigma(\mathfrak{R})$, we need to find a well-defined measure function.\par 
There are 2 ways to prove it's well-defined.
\begin{enumerate}
	\item \textbf{Young}
	\item \textbf{Caratheodory}
\end{enumerate} 

\begin{definition}[outer-measure]
	A outer measure $\eta: \mathfrak{B}(X) \rightarrow \bar{\mathbb{R}}$ is a map with the following properties:
	\begin{enumerate}
		\item $\eta(\emptyset) = 0$
		\item For $A \subset B$: $\eta(A) \leq \eta(B)$
		\item For a sequence $(A_n)_{n \geq 1}$: $\eta(\bigcup_n A_n) \leq \sum_{n = 1}^\infty \eta(A_n)$
	\end{enumerate}
\end{definition}

\begin{definition}
	Given a outer measure $\eta$, a set $A \subset X$ is $\eta$-measurable if:
	\begin{align*}
		\forall Q \subset X: \eta(Q) = \eta(Q \cap A) + \eta(Q \cap A^c)
	\end{align*}
\end{definition}

The convenience of Caratheodory is partially because the outer-measure defines a canonical $\sigma$-field. We can further prove the existence of an unique measure on the $\sigma$-field.
\begin{theorem}[Caratheodory]
	Given a outer measure $\eta: \mathfrak{B}(X) \rightarrow \mathbb{R}$:
	\begin{align*}
		\mathfrak{A}:= \{ A \in \mathfrak{B}(X): A \quad \eta-measurable \}
	\end{align*}
	defines a $\sigma$ algebra and $\eta|_{\mathfrak{A}}$ is a measure.
\end{theorem}

\begin{theorem}[The 1. extension theorem]
	Given $\mu: \mathfrak{R} \rightarrow \widehat{\mathbb{R}}$ a content and for $A$ define:
	\begin{align*}
		\eta(A):= \inf\{\sum_{n = 1}^\infty \mu(A_n): A_n \in \mathfrak{R}, \cup A_n \subset A \}
	\end{align*}
	Then $\eta$ is an outer measure.
\end{theorem}

\subsubsection*{Uniqueness of the Extension}
one of the central result is the following theorem, which is often used to prove the uniqueness of a random variable.
\begin{theorem}[Uniqueness of the extension]
	Given measure $\mu, \nu$ on the measurable space $(X, \mathfrak{A})$. $\mathfrak{E}$ a \textbf{intersection-stable }generator that satisfies:
	\begin{enumerate}
		\item $\forall E \in \mathfrak{E}: \mu(E) = \nu(E)$
		\item $\exists (E_n)_{n\geq 1} < \infty$ in $\mathfrak{E}$: $\bigcup_{n = 1}^\infty E_n = X$($\sigma$-finite)
	\end{enumerate}
\end{theorem}

\subsubsection*{Completeness of the Extension}

\begin{theorem}[continuity of measure]
	Suppose $\mu : \mathfrak{R} \rightarrow \mathbb{R}$ is a content defined on the ring $\mathfrak{R}$. Then the followings are equivalent:
	\begin{itemize}
		\item $\mu$ is a premeasure
	\end{itemize}
\end{theorem}

\begin{theorem}[Caratheodory Extension theorem]
\end{theorem}

Generator: Why do we always use the generator?\\
What is the $\sigma$-field defined by a random variable? -It is the domain.\\

Function class: continuous function with compact support\\
bounded linear operator\\
$L^1$ function\\
$L^p$ function\\

If we can find a continuous function that is not measurable(using Cantor function), then it is not in $L^1$, then we have a continuous function that is not integrable.

\begin{example}[Find a Lebegue measurable set that is not Borel measurable]
\end{example}

- Using Cantor function we can construct a function, where $\mu (g(C)) = 1$, thus $g^{-1}$ is a continuous but not measurable function. \\
The difference between Lebegue measurable function and Borel measurable function(The definition of the Lebegue measure. Why is the counter-example constructed using Cantor function is Lebegue measurable but not $\mathcal{L}-\mathcal{L}$ measurable.\\
Extension:
1. Extension of a field\\
2. Extension of the pre-measure\\
3. Extension from an outer measure to a measure.\\
3. Extension from Borel measurable function to Lebegue measurable\\ function.
3. Caratheodory\\



\subsection{The Lebesgue measure}
\begin{theorem}[Approximation theorem]
\end{theorem}



\section{Measurable and Integrable Functions}
Why is the Riemann integral not working for functions that are discontinuous on a set with positive measure.

\begin{theorem}[Young's Construction]
	For a non-negative function $f$ on $X$, f is measurable iff there exists a sequence of step functions $u_n$ with $u_n \uparrow f$. 
\end{theorem}
\begin{proof}
\end{proof}


\subsection*{numerical measurable functions}
\subsubsection*{Approximation by the simple functions}


\subsubsection*{countably generated measure space} 
First we could like to introduce the following definition:\\
1. \textbf{measurable isomorphism}:  \\
2. \textbf{$\mathfrak{G}$ separated $X$}: $\forall x, y \in X, x\neq y \exists A \in \mathfrak{G}: \chi_A(x) \neq \chi_A(y)$    \\
3. \textbf{$X$ is separated}: The $\sigma$-field $\mathfrak{A}$ separate $X$.\\
A space $(X, \mathfrak{A})$ is \textbf{countably generated}, if $\mathfrak{A}$ could be generated by a countable generator $\mathfrak{G}$: $\mathfrak{A} = \sigma{(\mathfrak{G}})$

\begin{lemma}
	$\mathfrak{G} \subset \mathfrak{A}$ separate $X$ iff $(X, \sigma(\mathfrak{G}))$ is separated
\end{lemma}

\begin{theorem}
	$(X, \mathfrak{A})$ is \textbf{countably generated} iff there exists a bijection $f: (X, \mathfrak{A}) \rightarrow ([0, 1], \mathfrak{B}([0, 1]))$.
\end{theorem}
\begin{proof}
	"$\Leftarrow$" is clear. For "$\Rightarrow$", see the function 
	\begin{align*}
		f:= \sum_{n = 1}^\infty 3^{-n} \chi_{A_n}
	\end{align*}
\end{proof}


\subsection{integrable function with compact support}
The functional spaces:
\begin{enumerate}
	\item $\mathcal{T}$: the step functions(simple functions). Have the form $f = \sum_{i = 1}^n \alpha \mu(A_i)$. 
\end{enumerate}
\begin{lemma}[the monotone convergence theorem]
	
\end{lemma}

\begin{theorem}
	$C_c(\mathbb{R}^n)$ is dense in $L^1(\mathbb{R}^n)$.
\end{theorem}
\begin{remark}
The technique of measure theory induction is applied here. Notice the result could be extended to $L^p$ space. See the chapter of separability.	
\end{remark}

\begin{corollary}
	\begin{align*}
		\forall \epsilon > 0 \exists \delta > 0 \forall E \in \mathfrak{A} \textbf{ with } \mu(E) < \delta:
		\int_E f d\mu < \epsilon
	\end{align*}
\end{corollary}


\subsection{integration of non-negative functions}

\subsection{The theorem of convergence}



\begin{theorem}[The dominated convergence theorem]
	Given measurable function $f_n, f : X \rightarrow \mathbb{K}$ with $f_n \overset{a.s.}{\rightarrow} f$. And $|f| \leq g$ a.s.  for some integrable function $g$. Then all $f_n$ and $f$ are integrable with:
	\begin{align*}
		\lim_{n \rightarrow \infty} \int f_n d\mu = \int f d\mu
	\end{align*}
	and 
	\begin{align*}
		\lim_{n \rightarrow \infty} \int |f_n - f| d\mu = 0
	\end{align*}
\end{theorem}
\begin{remark}
It's not enough to have $\mathbb{E}f_n < \infty$ for all $n$	
\end{remark}

\begin{example}[absolute continuous function]
	$f:[a, b] \rightarrow \mathbb{R}$ is differentiable and $f'$ is bounded. Then $f'$ is Lebesgue integrable over $[a, b]$ and 
	\begin{align*}
		\int_a^b f'd\lambda = f(a) - f(b)
	\end{align*}
\end{example}

\begin{theorem}[Fatu's lemma]
	For a sequence of measurable functions $f_n$ that non-negative, we have
	\begin{align*}
		\int \liminf f_n d\mu \le \liminf \int f_n d\mu 
	\end{align*}
\end{theorem}
\begin{proof}
\begin{align*}
	\int \lim_n \inf_{k \le n} f_k d\mu = \lim_n \int \inf_{k \le n} f_k d\mu \le \liminf \int f_n d\mu 
\end{align*}
\end{proof}

\begin{corollary}
	
\end{corollary}
\begin{proof}
\begin{align*}
	\int \lim_n \sup_{k \le k} f_k  = \lim_n \int \sup_{k \le n}f_k 
\end{align*}
\end{proof}



\newpage
\section{Product Measures}
The motivation of product measure is to define a measure in the cartesian product space which could be given by $\mu(A \times B) = \mu_X(A) \mu_Y(B)$. We want to explore the difference between 


\subsection{the product topology and the product measure}
\begin{theorem}
	A map $f: (Z, \mathfrak{G}) \rightarrow (X, \mathfrak{I})$ (where $\mathfrak{I}$ is the product $\sigma$-Algebra on the product space $X := \prod X_i$) is measurable iff it's measurable in every dimension.
\end{theorem}
\begin{proof}
\end{proof}
\quad\\
In the spirit of the last theorem we define the measurable function $j_K: \prod_{i \in I} X_i \rightarrow \prod_{i \in I} X_i$ :
\begin{equation*}
j_K(x) := \left\{
\begin{aligned}
	pr_i \qquad &i \in K\\
	a_i \qquad  &i \in I \backslash K \quad \textbf{for some $a_i \in X_i$}
\end{aligned}
\right.	
\end{equation*}
For any measurable set $M = \prod M_i \in \mathfrak{I}$ and $m_i \in M_i$ with respect to $a_i$, we can define the \textbf{cut} of $M$ as
\begin{align*}
	j_{I \backslash i} ^{-1}(\prod_{I \backslash i} \times a_i)
\end{align*}
By the last theorem we've established that every cut is measurable, therefore the following corollary
\begin{corollary}
	Every cut of a measurable set in the product $\sigma$-algebra is measurable.
\end{corollary}


\subsubsection{The measurability of the diagonal}
\begin{theorem}
	For the measurable space $(X, \mathfrak{A})$, the following claims are equivalent.
	\begin{enumerate}
		\item $X$ is separated by a countable $\mathfrak{G} \subset \mathfrak{A}$.
		\item There exists a bijection $f: (X, \mathfrak{A}) \rightarrow ([0, 1], \mathfrak{B}([0, 1]))$
		\item $\Delta_X \in \mathfrak{A} \otimes \mathfrak{A}$ 
		\item $\forall x \in X: \{x\} \in \mathfrak{A}$ and $(X, \mathfrak{A})$ is countably generated.
	\end{enumerate}
\end{theorem}
\begin{proof}
1)$\Rightarrow$2) trivial. \\
	3)$\Rightarrow$4): Notice: $\mathfrak{A} = \sigma(\mathfrak{G}) \Leftrightarrow \exists (A_n)_{n \geq 1} \in \mathfrak{G}: X = \bigcup_{n = 1}^\infty A_n$.(*)\\ 
	$\Delta_X \in \mathfrak{A} \bigotimes \mathfrak{A}$\\
	 $\Rightarrow \Delta_X \in \sigma(\mathfrak{G})$ with $\mathfrak{G} \subset \mathfrak{A} \otimes \mathfrak{A}$, countable.\\
	$\Rightarrow$ By (*) $\mathfrak{A} \otimes \mathfrak{A} = \sigma (\mathfrak{G})$\\
	$\Rightarrow$ $\mathfrak{A}$ is countably generated.
\end{proof}


\begin{definition}[Initial $\sigma$-Algebra and Product $\sigma$-Algebra]
	We have:
	\begin{itemize}
		\item\textbf{The initial $\sigma$-Algebra: } The smallest $\sigma$-Algebra on A, such that the family of function $\{f_i : A \rightarrow  (B, \mathfrak{B})| i \in I\}$ is measurable:
			\begin{equation*}
				\mathfrak{A} = \mathfrak{I}(f_i: i \in I) := \sigma(\cup f_i^{-1}(B_i) )
			\end{equation*} 
		\item \textbf{The product $\sigma$-Algebra}:
			The smallest $\sigma$-Algebra on A, such that the family of projection is measurable:
			\begin{equation*}
				\mathfrak{I}(pr_i: i \in I) := \bigotimes_{i \in I} \mathfrak{A}_i , \qquad pr_i: \prod X_j \rightarrow X_i, \quad  pr_i(x) = x_i
			\end{equation*}
		\item \textbf{Product Topology:} The smallest topology, such that the family of projection is continuous, a short deduction gives:
			\begin{equation*}
				(\prod X_j, \mathcal{T}) = \{ \prod_{\alpha \in I <\infty}  U_\alpha \times \prod_{\beta \neq \alpha} U_\beta \}
			\end{equation*}
	\end{itemize}
\end{definition}

\begin{remark}
Recall a \textbf{topology} is defined as a collection of subsets of: 1. contain $X$ and $\emptyset$, 2.any union, 3. pairwise intersection. 
 $B$ is a \textbf{subbase} if any subset in the topology is the union of the finite intersections of the subsets in the subbase. \par 
Given a space $X$ and a collection of subsets $B$ with $\cup_{U \in B} U = X$, the smallest topology that contains $B$ is the topological space generated by the subbase $B$.	
\end{remark}

\begin{theorem}
	$(X, \mathcal{T})$ is the \textbf{countable} product topology of the \textbf{separable} topological space $(X_i, \mathcal{T}_i$, then we have:
	\begin{equation*}
		\mathcal{B}(X) = \otimes_i \mathcal{B}(X_i)
	\end{equation*}
\end{theorem}

\begin{theorem}[Existence of Product measure]
$(X_j, \mathcal{A}_j, \mu_j)$ $(j = 1, \dots, n)$ \textbf{$\sigma$-finite measure space}, then there exists a \textbf{unique} measure:
\begin{equation*}
	\otimes_j \mu_j : \otimes \mathcal{A}_j \rightarrow \mathbb{R}^+, 
	\quad \otimes_j^n \mu_j(A_1 \times \dots \times A_n) = \prod_j^n \mu_j(A_j)
\end{equation*}
\end{theorem}


\subsection{Fubini theorem and the transformation rule}


\begin{theorem}[Fubini theorem]
	Given 2 \textbf{$\sigma$-finite} measure $\mu, \nu$ and define the canonical product measure $\mu \otimes \nu$ on the measurable space $(X \times B, \mathfrak{A} \otimes \mathfrak{B}$), given $f \in (A \times B, \mathfrak{A} \otimes \mathfrak{B}, \mu \otimes \nu)$ if one of the following criteria holds:
	\begin{enumerate}
		\item $f(x, \cdot)$ and $f(\cdot, y)$ non-negative.
		\item $f(\cdot, y)$ is integrable w.r.t $x$ for every $y$, and $f(x, \cdot)$ for every $x$.
		\item For $f$ measurable w.r.t the product $\sigma$-Algebra and either of the Integral is finite.
	\end{enumerate}
	Then we can change the order of integration.
\end{theorem}

\begin{theorem}[The Transformation Formula]
	Given $X, Y \subset \mathbb{R}^p$ open and $t: X \rightarrow Y$ a $C^1$-diffeomorphismus 
	\begin{enumerate}
		\item For all $A \in \mathcal{B}^p_X$, 
			\begin{equation*}
				\beta^p(t(A)) = \int_A |det Dt| d\beta^p
			\end{equation*}
	\end{enumerate}
\end{theorem}



\newpage
\section{Convergence Concept in the measure theory}
Throughout we assume the existence of $L^p$ space. The proof see Riezs-Fischer.


\subsection{Young, Holder and Minkowski}
\begin{proposition}[continuity of the convex functions on an Interval] Given a interval $I$, the convex function is continuous in the inner interval
\end{proposition}
\begin{proof}
	Page 241 of the Elstrodt.
\end{proof}

\begin{lemma}[Jensen's Inequality]
\end{lemma}
\begin{proof}
	The affine approximation and convex approximation.
\end{proof}

\begin{lemma}[Young's Ineqaulity]
if $1 < p, q < \infty$, $\frac{1}{p} + \frac{1}{q} = 1$, then:
\begin{equation*}
	ab \leq \frac{1}{p}a^p + \frac{1}{q}b^q
\end{equation*}
\end{lemma}

\begin{example}[An application of Jensen: Arithmetical and geometrical mean]
For $\sum_{i = 1}^n \alpha_i = 1$, $x_k > 0$, we have
	\begin{align*}
		\prod x_i^{\alpha_i} \leq \sum \alpha_i x_i 
	\end{align*}
\end{example}
\begin{proof}
	We see $\{ X := \{x_1, \dots, x_n\}, 2^X, \mu(x_i) = \alpha_i  \}$ defines a space with counting measure, For all measurable $Y: X \rightarrow\mathbb{R}$,  we therefore have by Jensen:
	\begin{align*}
		\mathbb{E}_X \exp(Y) \geq \exp(\mathbb{E} Y)
	\end{align*}
	choose $Y: y = \log x$, then
	\begin{align*}
		\sum_1^n \alpha_i x_i \geq \prod_1^n x_i^{\alpha_1} 
	\end{align*}
\end{proof}

\begin{lemma}[Holder's inequality]
$1 \leq p, q \leq \infty$, $\frac{1}{p} + \frac{1}{q} = 1$, $f, g$ measurable functions on $X$, we have:
\begin{align*}
	\|fg\|_{L^1} \leq \|f\|_{L^p} \|g\|_{L^q}
\end{align*}
\end{lemma}
\begin{proof}
	As we know from the results above:
	\begin{align*}
		\xi \eta \leq \frac{1}{p}\xi^p + \frac{1}{q}\eta^p
	\end{align*}
	Set $\xi = \frac{|f|}{\|f\|_p}, \eta = \frac{|g|}{\|g\|_q}$ and get the desired result.
\end{proof}
\begin{remark}
To apply the Holder usually we just need to make sure that $p > 1$.	
\end{remark}


\begin{theorem}[Minkowski Inequality]
	For $f, g \in L^1$ and $p \in [1, \infty]$, we have:
	\begin{align*}
		\| f + g \|_p \leq \|f\|_p + \|g\|_p
	\end{align*}
\end{theorem}
\begin{proof}
	\begin{align*}
		\|f + g\|^p_p =
		\int_\Omega |f + g|^p d\mu &= \int_\Omega |f+g|^{p-1} |f + g| d\mu \\
		&\leq \int_\Omega |f| |f+g|^{p-1} + |g| |f+g|^{p-1} d\mu \\
		&\leq (\|f\|_p + \|g\|_p) \|(f + g)^{p-1}\|_q\\
		&= (\|f\|_p + \|g\|_p) \|f + g\|_p^{\frac{p}{q}}\\
		(1)\Rightarrow
		\|f + g\|_p &\leq \|f\|_p + \|g\|_p
	\end{align*}
	In the case where division by infinity is not defined, we need to prove that $\|f + g\|_p$ is finite in (1).
\end{proof}

\begin{theorem}
	For $1 \leq p \leq \infty$ is $L^P$ a normed vector space and for $0 < p < 1$ is 
	\begin{equation*}
		d_p(f, g) := \|f - g\|^p_p
	\end{equation*}
	a metric in $L^p$
\end{theorem}


\subsection{The $L^p$ space}
Till now we constructed the basic topological structure of the $\mathcal{L}^p$ space. We established it's a vector space equipped with a metric(or in the case $p \geq 1$ a norm). Naturally we would like the space to be complete. The tools at hand are about pointwise convergence, with these tool we are able to deduce the $\mathcal{L}^p$ convergence.  

\begin{theorem}[Theorem of Riesz-Fisher]
	For $0 < p \leq \infty$ is $L^p$ a complete metric space w.r.t. $\|\cdot\|_p$. 
\end{theorem}
\begin{proof}
	The proof follows generally:
	\begin{itemize}
		\item Construct a pointwise convergent limit:
			\begin{enumerate}
				\item Use telescope sum to construct a subsequence in the form: $f_{n_k} = f_{n_1} + \sum g_i$. By the $L^p$ convergence of $f_n$ we know there exists a subsequence of $f_n$ such that 
				\begin{align*}
					\|f_{n_k} - f_{n_{k+1}}\|_p \leq 2^k\quad \forall m > n_k				\end{align*}
				define:
				\begin{align*}
					g_k := f_{n_k} - f_{n_{k+1}}
				\end{align*}
				\item $\lim_{k \rightarrow \infty} \sum_{i = 1}^k | g_i| := \hat{f}$ exists a.s. by bounded convergence theory since
				\begin{align*}
					 \|\sum |g_k|\|_p \le \sum \|g_k\|_p \le 1
				\end{align*}
				we have
				\begin{align*}
					\lim \|\sum|g_k|\|_p = \|\lim \sum |g_k|\|_p
				\end{align*}
				the pointwise limit is then well-defined. 
				\item \textbf{Construct the potential pointwise limit} 
					\begin{align*}
						\sum g_k = f_{n_1} - f_{n_k} \overset{a.s.}{\longrightarrow}  \hat{f}
					\end{align*}
					we define our potential pointwise limit
					\begin{align*}
						f:=  f_{n_1} - \hat{f}
					\end{align*}
			\end{enumerate}
		\item Show the convergence: At this point we have a a.s limit of the $L_p$ Cauchy sequence of function. We do not know if the limit is in $L^p$.
		\begin{enumerate}
			\item $\|f_n - f \|_p \rightarrow 0$
			\begin{align*}
				\int |f - f_n|^p d\mu = \int \liminf_k |f_{n_k} - f_m|^p d\mu \leq \liminf \int |f_{n_k} - f_m|^p d\mu 
			\end{align*}
			Since the sequence and the subsequence are $L^p$ convergent, we have that: $\liminf_k \int |f_{n_k} - f_m|^p d\mu \rightarrow 0$
			\item $f \in \mathcal{L}^p$
			\begin{align*}
				\int |f|^p d\mu = \int |f -f_m + f_m|^p d\mu 
				\leq 2^p \int |f-f_m|^p d\mu + 2^p \int |f_m|^p d\mu < \infty
			\end{align*}
		\end{enumerate}
	\end{itemize}
	The monotone convergence theory could be used to ensure the existence of a a.s. limit.
\end{proof}
	
\begin{corollary}[H.Weyl]
	For $0 < p \leq \infty$:
	\begin{enumerate}
		\item Every Cauchy sequence in $\mathcal{L}^p$ admits a a.s. convergent subsequence to $ f \in \mathcal{L}^p$
		\item Every convergent sequence in $\mathcal{L}^p$ admits a subsequence that converges a.s to a $f \in \mathcal{L}^p$.
	\end{enumerate}
\end{corollary}
	
\begin{corollary}
	For $1 \leq p \leq \infty$, $L^p$ is a Banach space and for $p \in (0, 1)$, $L^p$ is a complete metric space
\end{corollary}
\begin{remark}
Notice the for $p \in (0, 1)$, the norm function no longer satisfies the triangular inequality since it's not a convex function.	However $\|f\|_p^p$ is still a metric.
\end{remark}

\begin{example}[$L^p$ convergent yet not pointwise convergent sequence]
\end{example}

\begin{example}[Compactness of $L^p$ space]
	the $L^p$ space is in general not compact, in the sense that every $L^p$ bounded sequence needs not to admit a convergent subsequence.
\end{example}

\begin{theorem}
	Given $\mu(X) < \infty$, $0 < p < p' \leq \infty$, $L^{p'} \subset L^p$ and 
	\begin{align*}
		\|f\|_p \leq C(X, p, p') \|f\|_{p'}
	\end{align*}
	In another word, there exist a continuous embedding $L^{p'} \hookrightarrow L^p$.
\end{theorem}
\begin{proof}
	Write $\int_{M}|f|^p dm $ as $\int_M |f|^p \cdot 1 dm $ and apply the Holder's inequality. 
\end{proof}
\begin{remark}[Counterexample: $\mu(X) = \infty$]
Observe the sequence space $l^p(\mathbb{R}) \subset l^q(\mathbb{R})$ $\forall p \le q$. 
\end{remark}


\subsection{Convergence in measure space }
\begin{definition}[Almost everywhere convergence]
Let $f, f_n : X \to \mathbb{K}$ ($n \in \mathbb{N}$) be measurable. The sequence $(f_n)$ \textbf{converges to $f$ almost everywhere} (written $f_n \to f$ $\mu$-a.e.) if there exists a null set $N \in \mathfrak{A}$ such that
\begin{equation*}
	f_n(x) \to f(x) \quad \text{for every } x \in X \setminus N.
\end{equation*}
\end{definition}

\begin{definition}[Uniform convergence a.e.\ and almost uniform convergence]
Let $f, f_n : X \to \mathbb{K}$ ($n \in \mathbb{N}$) be measurable. We distinguish two notions of uniform convergence on a measure space:
\begin{itemize}
	\item $(f_n)$ \textbf{converges uniformly almost everywhere} to $f$ if there exists a null set $N \in \mathfrak{A}$ such that
	\begin{equation*}
		\|f_n - f\|_{L^\infty(X \setminus N)} \xrightarrow[n \to \infty]{} 0,
	\end{equation*}
	equivalently, $\|f_n - f\|_{L^\infty(X)} \to 0$ (convergence in $L^\infty$).
	\item $(f_n)$ \textbf{converges almost uniformly} to $f$ if for every $\delta > 0$ there exists a measurable set $E \in \mathfrak{A}$ with $\mu(E) < \delta$ such that
	\begin{equation*}
		\|f_n - f\|_{L^\infty(X \setminus E)} \xrightarrow[n \to \infty]{} 0,
	\end{equation*}
	i.e.\ $(f_n)$ converges to $f$ uniformly on the complement $X \setminus E$.
\end{itemize}
The key difference is that in uniform convergence a.e.\ the exceptional set $N$ is fixed (and of measure zero), whereas in almost uniform convergence the exceptional set $E$ may depend on $\delta$ (and is only required to have arbitrarily small measure, not zero).
\end{definition}
\begin{remark}
Notice a.s. uniform convergence is strictly stronger, since the exceptional null set $N$ is fixed once and for all, whereas in almost uniform convergence the exceptional set $E$ is allowed to depend on $\delta$. The canonical witness is $f_n(x) = x^n$ on $[0, 1]$:
\begin{itemize}
	\item $f_n \to 0$ pointwise on $[0, 1)$ and $f_n(1) = 1$, so $f_n \to 0$ a.e.
	\item $f_n \not\to 0$ a.s. uniformly: for any null set $N \subset [0, 1]$ the set $[0, 1) \setminus N$ still has supremum $1$, so $\|f_n\|_\infty$ on $[0, 1] \setminus N$ does not tend to $0$.
	\item $f_n \to 0$ almost uniformly: for any $\delta > 0$ take $E = (1 - \delta, 1]$ with $\mu(E) = \delta$; on $[0, 1] \setminus E = [0, 1 - \delta]$ we have $\sup x^n = (1 - \delta)^n \to 0$.
\end{itemize}
The dependence of the exceptional set on $\delta$ is precisely what allows the singularity at $x = 1$ to be cut out at arbitrarily small cost.
\end{remark}

\begin{definition}[Convergence in measure]
Let $f, f_n : X \to \mathbb{K}$ ($n \in \mathbb{N}$) be measurable functions on a measure space $(X, \mathfrak{A}, \mu)$.
\begin{itemize}
	\item The sequence $(f_n)_{n \geq 1}$ \textbf{converges in measure} to $f$, written $f_n \xrightarrow{\mu} f$, if for every $\varepsilon > 0$,
	\begin{equation*}
		\lim_{n \to \infty} \mu\bigl(\{|f_n - f| \geq \varepsilon\}\bigr) = 0.
	\end{equation*}
	\item The sequence $(f_n)_{n \geq 1}$ \textbf{converges locally in measure} to $f$, written $f_n \xrightarrow{\mathrm{loc}\,\mu} f$, if for every $\varepsilon > 0$ and every $A \in \mathfrak{A}$ with $\mu(A) < \infty$,
	\begin{equation*}
		\lim_{n \to \infty} \mu\bigl(\{|f_n - f| \geq \varepsilon\} \cap A\bigr) = 0.
	\end{equation*}
\end{itemize}
\end{definition}

\begin{remark}[Ky Fan metric]
On a finite measure space ($\mu(X) < \infty$), convergence in measure is metrizable. For measurable $f, g : X \to \mathbb{R}$ set
\begin{equation*}
	\delta(f, g) := \inf \bigl\{ \varepsilon \geq 0 : \mu\bigl(\{|f - g| > \varepsilon\}\bigr) \leq \varepsilon \bigr\}.
\end{equation*}
The infimum is attained, $\delta$ is a semi-metric on the space $\mathcal{M}(X, \mathbb{R})$ of measurable functions, and $f_n \xrightarrow{\mu} f$ iff $\delta(f_n, f) \to 0$. In particular, the notion of Cauchy sequence in measure is well defined. Equivalently, the semi-metric
\begin{equation*}
	d(f, g) := \int_X \frac{|f - g|}{1 + |f - g|} d\mu
\end{equation*}
metrizes convergence in measure on a finite measure space.
\end{remark}

\begin{remark}
In probability theory, convergence in measure is called \emph{convergence in probability}.
\end{remark}


We'll now explore the relation between those convergence concept. First we study the relation between almost uniform and the a.s pointwise convergence. The conclusion is: almost uniform convergence implies a.s. pointwise convergence. The converse is true if we have a finite measure space. This is proved by the Russian mathematician \textbf{Jegorow}.

\begin{lemma}
	Given measurable function $f_n, f : X \rightarrow \mathbb{R}$, we have:
	\begin{itemize}
		\item $f_n \overset{a.s.}{\rightarrow} f \Leftrightarrow \forall \epsilon > 0:
		  \mu( \{x:   \bigcap_{n = 1}^\infty  \bigcup_{k = 1}^\infty |f_{n + k}(x) - f(x)| > \epsilon \}) = 0 $
		  \item $f_n \overset{a.s.}{\rightarrow} f \Leftarrow \forall \epsilon > 0 : \lim_{n \rightarrow \infty} \mu(\{x: \bigcup_{k = 1}^\infty | f_{n+k}(x) - f(x)| \geq \epsilon \}) = 0 $
		  \item $f_n \overset{a.s.}{\rightarrow} f \Leftrightarrow \forall \epsilon > 0  : \lim_{n \rightarrow \infty} \mu(\{x: \bigcup_{k = 1}^\infty | f_{n+k}(x) - f(x)| \geq \epsilon \}) = 0 $ if $\mu(X) < \infty$.		  
	\end{itemize}
\end{lemma}
\begin{proof}
	a) is a reformulation of the definition of convergence, b) follows directly form a) and c) follows form b) and the continuity of measure from above. 
\end{proof}
\begin{remark}
The equivalent definition of a.s. convergence could be applied to Cauchy-sequence as well.	
\end{remark}

\begin{lemma}
	Every \textbf{almost uniformly convergent} sequence $f_n$ admits an \textbf{a.s pointwise convergent} limit $f \in L^0$.
\end{lemma}
\begin{proof}
	Define $E_k$ with $\mu(E_k) < \frac{1}{k}$ and $f_n \overset{\infty}{\rightarrow} f$ on $X \slash E_k$. Observe $E:= \bigcap_{k = 1}^\infty E_k$\\ 
	The existence of $E$ is given by the fact that we have a $\sigma$-field. And 
	\begin{align*}
		\mu(E) = \lim_{k \rightarrow \infty} \mu (\bigcap_{k = 1}^{\infty} E_k) \leq \lim_{k \rightarrow \infty} \frac{1}{k} \rightarrow 0
	\end{align*}
	We can argue by contradiction that $\forall x \in E^c$, the function converges.
\end{proof}
\begin{remark}
From \textbf{uniform convergence on $E_n$} we cannot follow \textbf{uniform convergence on $\bigcup^\infty_1 E_n$}. However, we get \textbf{pointwise convergence on $\bigcup_1^\infty E_n$}.	
\end{remark}

\begin{theorem}[the Jegorow's theorem]
	Given a finite measure space($\mu(X) < \infty$),  \textbf{a.s pointwise convergence} implies \textbf{almost uniform convergence}.
\end{theorem}
\begin{proof}
	Uniform
\end{proof}

\begin{theorem}
	Almost uniform convergence implies convergence i.M.
\end{theorem}

\begin{theorem}
	Every i.M convergent sequence admits a almost uniform convergent subsequence
\end{theorem}
\begin{proof}
	By convergence in measure we have: $\forall n,m \geq n_k$
	\begin{align*}
		\mu(|f_m - f_n| \geq 2^{-k}) \le 2^{-k}
	\end{align*}
	Define:
	\begin{align*}
		A_k := \{ x \in \Omega| f_m - f_n \geq 2^{-k}\}
	\end{align*}
	we have 
	\begin{align*}
		B_l := \cup_{k = l}^\infty A_k ,\qquad 
		\mu(B_l) \le \sum_{k = l}^\infty 2^{-k} = 2^{-l}
	\end{align*}
	Then we have $\forall \mu(B_m^C) \le \delta), \forall k > l > m$:
	\begin{align*}
		|f_{n_k} - f_{n_l}| \le \sum_{i = l+ 1}^k |f_{i} - f_{i -1}| \le 2^{-m}
	\end{align*}
\end{proof}

\begin{theorem}
	Given a $\sigma$-finite measure space, every locally i.m convergent sequence admits an a.s. pointwise convergent subsequence.
\end{theorem}
\begin{proof}	
\end{proof}


\subsection{Convergence in \texorpdfstring{$L^p$}{Lp} space}
\begin{definition}[uniformly $p$-integrable]
A measurable sequence of functions $(f_n)$ is uniformly $p$-integrable if:
\begin{itemize}
	\item tightness: $\forall \epsilon > 0 \exists E \in \mathfrak{A} \textbf{ with } \mu(E) < \infty: \sup_n \int_{E^c} |f_n|^p d\mu < \epsilon$
	\item no concentration: $\forall \epsilon > 0 \exists \delta > 0 \forall E \in \mathfrak{A} \textbf{ with } \mu(E)< \delta: \sup_n \int_E |f_n|^p d\mu < \epsilon $ 
\end{itemize}
\end{definition}

The conditions could be summarized as:\par 
The functions are supported \textbf{arbitrarily well} on a finite-measure set. Therefore the integral cannot explode on a small set.\par 
For all the functions, the integral \textbf{is continuous at 0} with respect to the measure at 0. Therefore the measurable set supporting an arbitrary small integral cannot explode.

\begin{theorem}[The Vitali convergence theorem]
	For $0< p < \infty$ and $f_n, f$ measurable in $L^p$, if:
	\begin{itemize}
		\item $f_n \textbf{  u.i.}$
		\item $f_n \rightarrow f$ locally in measure
	\end{itemize}
	Then we have:
	\begin{align*}
		f_n \overset{L^p}{\longrightarrow} f 
	\end{align*}
\end{theorem}
\begin{proof}.\\
	"$\Rightarrow$":\\
	"$\Leftarrow$"\\
	First we prove if $f_n \overset{a.s}{\rightarrow} f$, the claim holds. 
	Select $E$ as in (1) of the definition of u.i. and $E \slash B$ as in the (2). We have:
	\begin{align*}
		\int_\Omega |f - f_n|^p d\mu \leq 2^p \int_{E^c} |f|^p + |f_n|^p d\mu 
		+ 2^{p-1} \int_{E \slash B} |f|^p + |f_n|^p d\mu + \int_B |f - f_n|^p d\mu 
	\end{align*}
	By Fatou:
	\begin{align*}
		\int_{E^c} |f|^p d\mu \leq \liminf \int_{E^c} |f_n|^p d\mu \leq \epsilon\\
		\int_{E\slash B} |f|^p d\mu  \leq \liminf \int_{E^c} |f_n|^p d\mu  \leq \epsilon
	\end{align*}
	Together with the definition of u.i., we conclude that the first terms approach to 0. And by the Jegorow's theorem we know on $B$ $f_n$ convergence uniformly therefore in $L^p$.\par 
	Now argue by contradiction, assuming 
	\begin{align*}
		\exists \epsilon >0 \forall k \in \mathbb{N} : \|f_{n_k} - f\|_p > \epsilon
	\end{align*}
	$f_{n_k}$ is a subsequence of a convergent(in measure) sequence therefore converges in measure. We get an a.s. pointwise convergent subsequence of $f_{n,k}$ and it convergence in $L^p$ to $f$, which contradicts the assumption.
\end{proof}
\begin{remark}
The technique $(f \pm g)^p \leq (\frac{2|f| + 2|g|}{2})^p \leq \frac{2^p |f|^p + 2^p|g|^p}{2} = 2^{p-1}(|f|^p + |g|^p)$ is often used in the context of $L^p$ convergence.	
\end{remark}



\newpage
\section{Absolute Continuity}
\begin{definition}[Absolute continuous function]
A function $F$ is absolute continuous if:$\forall \epsilon > 0 \exists \delta > 0$, such that for all finite partition $(a_i, b_i)$: 
\begin{equation*}
	\sum_{i = 1}^N b_i - a_i < \delta \Rightarrow \sum_{i = 1}^N |F(b_i) - F(a_i)| < \epsilon
\end{equation*}

\end{definition}
\begin{remark}
\end{remark}

\begin{definition}[Locally integrable function]
	\begin{equation*}
		L^p_{loc}(\Omega) := \{ f: \Omega \rightarrow V | f|_K \in L^p(\Omega) \forall K \subset \Omega \text{ compact} \}   \}
	\end{equation*}
\end{definition}

\begin{definition}[Lebesgue Point]
	For $f \in L_{loc}^1(\mathbb{R}^n), x \in \mathbb{R}$ is called a Lebesgue point of $f$ if:
	\begin{equation*}
		\lim_{r \rightarrow 0^+} \frac{1}{m(B_r(x))} \int_{B_r(x)} |f(x) - f(y)|dy = 0
	\end{equation*}
\end{definition}

\begin{definition}[Maximal Function]
	For $ f \in L^1_{loc}(\mathbb{R}^n)$ the \textbf{maximal function} $Mf: \mathbb{R}^n \rightarrow [0, \infty]$ is defined as:
	\begin{equation*}
		Mf(x) = \sup_{r > 0} \frac{1}{m(B_r(x))} \int_{B_r(x)} |f|dm
	\end{equation*}
\end{definition}

\begin{lemma}
	We define $f^0 := \limsup_{r \rightarrow 0^+} \frac{1}{m(B_r(x))} \int_{B_r(x)} |f|dm$.For every $N \in \mathbb{N}$, 
	\begin{equation*}
		A_N := \{ x \in \mathbb{R}^n | f^0(x) > \frac{1}{N} \}
	\end{equation*}
	Then $A_N$ is contained in a Lebesgue measurable set with measure less than $\frac{1}{N}$.
\end{lemma}

\begin{theorem}[Lebesgue Differentiation theorem]
	For $f \in L_{loc}^1(\mathbb{R}^n)$. Almost every point is a Lebesgue point.
\end{theorem}
\begin{remark}
To see why it is called the differentiation theorem, observe:
\begin{equation*}
	\frac{F(x + h) - F(x)}{h} - f(x) = \frac{1}{h}\int_x^h f(t) - f(x) dt \leq 2 \frac{1}{m(x-h, x + h)} \int^{x+h}_{x-h} |f(t) - f(x)| dt \rightarrow 0 
\end{equation*}	
\end{remark}


\begin{corollary}
	$\forall f \in L^1([a, b])$, $F := \int_a^b f(t) dt$ is differentiable a.e. and $F' = f$ 
\end{corollary}

\begin{theorem}[Fundamental theorem of Calculus for the Lebesgue Integral]
\end{theorem}

\begin{theorem}[Radon-Nikonym]
	Given $\mu$ a $\sigma$-finite measure and $\nu \ll \mu$ a signed measure on $\mathcal{A}$, then there exist a quasi-integrable function $f: X \rightarrow \bar{\mathbb{R}}$, such that $\nu = f \circ \mu$ . 
\end{theorem}
\begin{proof}
	w.l.o.g. we only proved the case of positive measure.
	\begin{itemize}
		\item \textbf{Prove the lemma:} If $\mu, \nu$ are finite measure and $\nu \leq \mu$, then there exists a $\rho$, such that $\nu = \rho \circ \mu$ with $\rho: X \rightarrow [0, 1]$. And $h$ is unique in $L^2$. \par
			 
		\begin{enumerate}
			\item For a partition $P_1$ of $X$, define the stepfunction:
				\begin{equation*}
					p_1(x) =\left\{ \begin{aligned}
						\frac{\nu(A)}{\mu(A)} \quad &\mu(A) \neq 0 \\
						0 \quad &\mu(A) = 0
					\end{aligned}
					\right.
				\end{equation*}
			Given a finer partition $P_2$, we know $\forall A \in P_1:$
				\begin{equation*}
					\int_A p_1 d\mu = \int_A p_2 d\mu 
				\end{equation*} 
			\item Define:
				\begin{equation*}
					\alpha = \sup\{\int_X p^2 d\rho |\text{ $P$ a partition of X} \}
				\end{equation*}
				and $p_n$ such that:
				\begin{equation*}
					\alpha \geq \int_X p_n^2 d\rho \geq \alpha - \frac{1}{n} 
				\end{equation*}
				then for the finest common partition:
				\begin{equation*}
					\int_X (p_n - p_m)^2 d \mu = \int_X p_n^2 - p_m^2 d\rho \leq \frac{1}{m}
				\end{equation*}
				which makes it a Cauchy sequence. Since $L^2$ is complete, we know $\|p_n - h\|_2 \rightarrow 0$ for a $h \in L^2$
		\end{enumerate}
		\item \textbf{Prove the case: $\nu$ and $\mu$ are finite.}\par 
		We define $\tau = \mu + \nu$, sine now $mu \leq \tau$ and $\nu \leq \tau$, we can apply the lemma above, and get $\tau = g \circ \mu$ and $\tau = h \circ \nu$. Now define $N:= \{M \subset X | \mu(M) = 0\}$ and $f$:
		\begin{equation*}
			f(x) = \left\{ \begin{aligned}
				\frac{h(x)}{g(x)} \quad & x \in N^c\\
				0 \quad & x \in N
			\end{aligned}\right.
		\end{equation*}
		then $\forall A \in \mathcal{A}$( $N$ excluded):
		\begin{equation*}
			\nu(A) = \int_A h d \tau = \int_A f \circ g d\tau = \int_A f d \mu
		\end{equation*}
		\item Prove the case: $\mu$ is finite, $\nu$ arbitrary. 
		\item Prove the case: $\mu$ $\sigma$-finite, $\nu$ arbitrary. 
	\end{itemize}
\end{proof}


\subsection{Hahn decomposition}
\begin{definition}[the signed measure]
	Given $\mathcal{A}$ the $\sigma$-algebra, the function $\nu: \mathcal{A} \rightarrow \bar{\mathbb{R}}$ is called the \textbf{signed measure} if 
	\begin{enumerate}
		\item $\nu(\emptyset) = 0$,
		\item $\nu(\mathcal{A}) \subset (-\infty, \infty]$ or $\nu(\mathcal{A} \subset [-\infty, \infty)$,
		\item $\sigma$-additivity. 
	\end{enumerate}
\end{definition}


\begin{theorem}[Hahn decomposition]
Let $\nu$ be a signed measure $(X, \mathcal{A}) \rightarrow \bar{\mathbb{R}}$. For every $\mu$ there exists a decomposition of $X$: $X = P \cup N$ ($P, N \in \mathcal{A}$) such that 
\begin{enumerate}
	\item $\forall A \in \mathcal{A}, A \subset P:\quad \nu(A) \ge 0$, i.e., $P$ is $\nu$-positive. 
	\item $\forall B \in \mathcal{A}, B \subset N: \quad \nu(B) \le 0$, i.e., N is $\nu$-negative. 
	\item $P$ and $N$ are unique outside a $\nu$-null set. 
\end{enumerate}
\end{theorem}




\section{Measure in topological spaces}
\begin{definition}
	Given a measurable space $(X, \mathcal{F})$, a measure $\mu: \mathcal{B} \rightarrow [0, \infty]$ is called 
	\begin{enumerate}
		\item \textbf{locally finite measure (Borel measure)}: 
			\begin{equation*}
				\forall x \in X \exists U \ni x: \mu(U) < \infty
			\end{equation*}
		\item \textbf{Inner \& outer regular measure}: $\mu$ is inner regular if:
			\begin{equation*}
				\forall A \in \mathcal{B}: \mu(B) = \sup_{K \in \mathcal{B} \text{compact}} \mu(K) 
			\end{equation*}
		\item \textbf{Radon-measure:} a inner regular Borel-measure.
		\item \textbf{tight}
	\end{enumerate}
\end{definition}

\begin{theorem}[The theorem of Ulam]
\end{theorem}

\begin{proof}
	\begin{itemize}
		\item \textbf{The lemma:} Every finite measure on a $\sigma$-Algebra, the set:
			\begin{equation*}
				\{A \subset \mathcal{A}: A \text{ is  $\mu$-regular}\}
			\end{equation*}
			is a $\sigma$-ring
	\end{itemize}
\end{proof}

\begin{definition}[$\mu$-boundry-less set]
In the measure space $(X, \mathcal{B}, \mu)$, a set $B$ is called $\mu$-boundary-less, if:
\begin{equation*}
	\mu(\partial B) = \mu(\bar{B} \setminus \overset{\circ}{B}) = 0
\end{equation*}
\end{definition}

\begin{theorem}[Portmanteau-Theorem]
	For a metric space $X$, $\mu_n, \mu$ finite Borel measure, then the followings are equivalent:
	\begin{enumerate}
		\item $\mu_n \rightharpoonup \mu$
		\item $\forall F \subset X \text{ closed}: \limsup \mu_n(F) \leq \mu(F)$
		\item $\forall U \subset X open: \liminf \mu_n(U) \geq \mu(F)$
	\end{enumerate}
\end{theorem}
\newpage



